%%% Local Variables:
%%% mode: latex
%%% TeX-master: "../main"
%%% End:

\chapter[绪论]{绪论}
\label{cha:intro}

\section{研究背景与意义}

\subsection{两轮足机器人概述}

足式机器人具有离散支撑、地形适应能力强的特点，能够在台阶、碎石、斜坡等非结构化地形中灵活运动\cite{raibert1986}；轮式机器人则凭借连续的滚动接触，具备高速、高效、低能耗与高平稳性的优势\cite{klemm2019}。两轮足机器人（Two-Wheeled Legged Robot）以两个主动驱动轮作为运动执行末端，通过腿部机构的姿态调节实现机身的平衡与越障，本质上是一种可主动平衡的欠驱动倒立摆系统，兼具两类移动方式的优点。

与常规四足或双足机器人相比，两轮足机器人结构更轻、自由度更少，其稳定站立依赖轮端力矩的实时平衡调节，而非足底静平衡，因此对控制器的动态响应与鲁棒性要求极高\cite{raibert1986}。近年来，两轮足机器人在安防巡检、仓储物流、楼宇服务等场景中展现出广阔的应用前景，围绕其开展的步态生成与控制研究成为移动机器人领域的热点方向。

\subsection{强化学习在机器人运动控制中的应用}

传统机器人运动控制通常依赖精确的系统动力学建模与模型预测控制（MPC）等方法，在模型误差大、环境复杂多变的场景下难以保证控制性能。随着深度强化学习（Deep Reinforcement Learning, DRL）的发展，策略网络能够在仿真环境中通过大量试错自主学习控制策略，已成为足式机器人步态生成的主流范式\cite{kober2013,hwangbo2019}。

具体而言，基于强化学习的运动控制在四足机器人上取得了显著进展：Rudin等\cite{rudin2022}利用大规模并行仿真在数分钟内完成四足机器人行走策略训练；Lee等\cite{lee2020}通过学习地形感知策略实现了四足机器人在复杂地形上的鲁棒运动；Miki等\cite{miki2022}进一步将感知与运动控制策略结合，实现了野外感知运动。在双足与人形机器人领域，强化学习同样被用于实现跳跃、空翻等高动态敏捷运动\cite{margolis2022,cheng2024}。这些成果表明，强化学习为两轮足机器人多步态控制提供了可行的技术路径。

\subsection{课题研究意义}

两轮足机器人若要完成"多地形、多步态、多算法"的综合任务，需要策略具备连续学习多种步态的能力。然而，当同一策略网络先后学习滚动行走、抬腿行走、跳跃、后空翻、抬腿上台阶与单腿平衡等多种步态时，会面临三个核心科学问题：一是\textbf{灾难性遗忘}，学习新步态可能导致旧步态能力退化；二是\textbf{奖励冲突}，如"跳得高"与"保持平衡"等子目标相互矛盾；三是\textbf{局部最优}，策略可能陷入"剧烈晃动"等取巧捷径。

本文以自主研制的两轮足机器人为研究对象，围绕上述三个核心问题，构建"渐进式-自适应-强鲁棒"的多步态运动控制与调度框架，系统研究多专家网络调度、渐进式课程学习、多算法融合与精细平衡控制等关键技术。研究成果对推动两轮足机器人工程化应用具有重要的理论价值与现实意义。

\section{国内外研究现状与分析}

\subsection{两轮足与轮腿机器人运动控制研究现状}

两轮足机器人本质上属于轮腿混合机器人（Wheel-Legged Robot）的一个子类。在机构设计方面，瑞士苏黎世联邦理工学院（ETH Zurich）研制的Ascento机器人采用四连杆腿部机构配合双轮驱动，实现了高速移动与跳跃运动\cite{klemm2019}。在国内，相关团队也开展了轮腿复合机构、可变形轮足机构的研究，通过关节复用实现足式模式与轮腿模式的切换。

在控制方法方面，早期研究主要采用PID、LQR等线性控制器对两轮自平衡系统进行镇定，其后基于SLIP（Spring-Loaded Inverted Pendulum）模型的柔顺控制与基于虚拟模型控制（Virtual Model Control, VMC）的力分配方法被引入轮腿机器人\cite{pratt2001,blickhan1989}。近年来，模型预测控制（MPC）与全状态反馈控制也被用于轮腿机器人的平衡与导航。但上述方法依赖精确动力学模型，面对复杂多变的地形与高动态步态时参数整定困难、泛化能力有限。

\subsection{基于强化学习的足式机器人步态生成研究现状}

强化学习为足式机器人步态生成提供了不依赖精确模型的控制框架。Hwangbo等\cite{hwangbo2019}提出通过仿真与真实域的联合训练，使四足机器人学习快速、动态的敏捷技能；Lee等\cite{lee2020}使用教师-学生框架与领域随机化，实现四足机器人对复杂地形的鲁棒适应；Rudin等\cite{rudin2022}证明了大规模并行仿真能够极大加速强化学习训练进程，使行走技能在数分钟内即可习得。

在敏捷运动方面，Margolis等\cite{margolis2022}利用强化学习使四足机器人学习跳跃技能，并通过仿真-现实迁移实现真实跳跃；Zhuang等\cite{zhuang2023}提出了机器人跑酷学习框架，将复杂敏捷技能分解为可组合的基本技能；Cheng等\cite{cheng2024}进一步实现了四足机器人的极限跑酷与后空翻。在轮腿机器人方面，已有研究将强化学习用于两轮自平衡机器人的平衡与路径跟踪\cite{klemm2019}，并逐步向复杂地形与敏捷动作拓展。

\subsection{连续多任务学习与灾难性遗忘研究现状}

当策略网络按顺序学习多个任务时，网络参数在新任务上的更新会覆盖旧任务所学的知识，导致灾难性遗忘（Catastrophic Forgetting）现象\cite{mccloskey1989}。为缓解该问题，主流方法包括：弹性权重巩固（EWC）等基于正则化的方法\cite{kirkpatrick2017}、基于记忆回放的持续学习方法、以及网络结构扩展方法。在机器人运动控制领域，采用渐进式课程学习（Curriculum Learning）从易到难组织训练任务，已被证明能够有效降低遗忘风险并提升训练稳定性\cite{bengio2009}。

此外，基于专家网络的混合结构（Mixture of Experts, MoE）通过门控机制动态组合多个子网络，天然适合多步态并存的控制场景\cite{jacobs1991,shazeer2017}。如何将课程学习与专家调度结合，在两轮足机器人连续多任务学习场景下缓解灾难性遗忘，仍是一个有待深入研究的问题。

\subsection{约束强化学习与融合算法研究现状}

两轮足机器人的高动态步态往往需要在安全性约束（如关节力矩极限、落地冲击、姿态范围）下优化性能指标。约束马尔可夫决策过程（CMDP）将约束引入目标函数，策略优化需在满足约束的前提下最大化累计回报\cite{cpo2017}。N3PO（Natural Policy Gradient with safety）等基于自然策略梯度的约束优化算法，通过拉格朗日对偶或投影方法在策略更新中显式满足约束，适合处理"既要高动态又要稳姿态"的矛盾目标\cite{npg2001}。

在算法融合方面，将模型先验与强化学习结合是提升样本效率与控制质量的常用思路：SLIP模型为跳跃等弹道运动提供解析的先验运动规划\cite{blickhan1989,geyer2006}；虚拟模型控制（VMC）为平衡与力分配提供直观的虚拟力-关节力矩映射\cite{pratt2001}。本文正是基于"PPO+SLIP""SLIP+PPO+VMC"等融合范式展开跳跃与后空翻控制研究。

\subsection{研究现状分析与存在的问题}

综上所述，国内外在足式机器人强化学习控制、敏捷运动生成与轮腿机器人机构设计等方面取得了丰硕成果，但面向两轮足机器人的研究仍存在以下不足：

（1）多步态连续学习场景下，将多种步态直接纳入单一策略训练会造成灾难性遗忘，而"各步态独立成专家、再统一调度整合"的"分训-整合"范式尚未在两轮足机器人上得到系统研究；

（2）跳跃、后空翻等高动态步态与平衡、上台阶等精细任务之间存在奖励冲突，单一算法难以兼顾，需要融合算法与约束优化的系统对比研究；

（3）多步态并存时如何依据地形与任务状态自动调度不同的专家网络，缺乏学习式、自适应的统一调度机制及其在混合地形环境中的整体验证。

上述问题构成了本文的研究出发点。

\section{本文主要研究内容}

本文以两轮足机器人为对象，围绕"分训-整合"的总体思路，构建由7个步态专家网络与1个学习式统一调度器组成的多任务强化学习框架，最终实现在混合地形环境中按地形自动切换步态。全文研究内容如下：

（1）\textbf{步态专家独立训练（分训阶段）}。针对7类典型步态分别设计并独立训练步态专家策略网络：滚动行走专家（含前进后退、转向与侧向移动）、粗糙地形行走专家、抬腿行走专家、跳跃专家、后空翻专家、抬腿上台阶专家与单腿平衡专家。各专家在对应地形的专用环境中训练，并针对不同算法范式（单PPO、PPO+SLIP、SLIP+PPO+VMC、N3PO）开展对比研究。

（2）\textbf{学习式统一调度器设计（整合阶段）}。设计以地形感知与本体状态为输入的学习式调度器，通过分层强化学习在混合地形统一环境中训练，实现机器人按地形自动选择与平滑切换步态专家：粗糙地形采用抬腿点足行走、低台阶抬腿跨越、高台阶跳跃上去、狭窄地形侧向移动、独木桥单腿平衡、木桩间隙跳跃通过。

（3）\textbf{统一系统集成与整体验证}。将7个步态专家与学习式调度器集成为完整系统，在由平面、粗糙、碎石、台阶、独木桥、木桩与狭窄通道构成的混合地形统一环境中进行整体遍历验证，并对多任务整合中的灾难性遗忘、奖励冲突与局部最优三大问题进行消融分析。

本文总体研究架构如图~\ref{fig:architecture}所示，涵盖多地形、多步态与多算法三个维度，各研究内容在此架构下按"分训-整合"范式展开。

\begin{figure}[!htbp]
    \centering
    \includegraphics[width=0.98\textwidth]{两轮足机器人架构.png}
    \bicaption{两轮足机器人总体研究架构（多地形-多步态-多算法）}{Overall research architecture of the two-wheeled legged robot (multi-terrain, multi-gait, multi-algorithm)}
    \label{fig:architecture}
\end{figure}

\section{论文组织结构}

本文共分九章，组织结构如图~\ref{fig:structure}所示。

\begin{figure}[!htbp]
    \centering
    
\includegraphics[width=0.95\textwidth]{WUST-flash.pdf}
    \bicaption{论文组织结构图（替换为自绘图）}{Thesis organization structure (replace with your own figure)}
    \label{fig:structure}
\end{figure}

第1章为绪论，阐述研究背景、国内外研究现状与主要研究内容；第2章介绍两轮足机器人建模与强化学习理论基础；第3章设计"分训-整合"的多任务强化学习总体框架与混合地形统一环境；第4-6章分别阐述行走类、精细平衡类与高动态敏捷类步态专家的独立训练；第7章设计学习式统一调度器；第8章进行统一系统集成与整体验证；第9章为结论与展望。
